\Subhead{1.9.2}{Разложение вынужденных колебаний по системе собственных функций}
\paragraph{Возбуждение каждой моды определяется проекцией источника.}
При однородных граничных условиях подставим $u=\sum_nQ_n(t)\phi_n$ в~\eqref{eq:distributedmodes}, умножим на $\phi_m^*$ и проинтегрируем. Благодаря ортогональности каждая нормальная координата подчиняется уравнению осциллятора:
\begin{equation}
 \ddot Q_n+2\gamma_n\dot Q_n+\omega_n^2Q_n=F_n(t),\qquad
 F_n(t)=\int_V\phi_n^*(\bm r)f(\bm r,t)\,\dd V.
\end{equation}
Слабые потери здесь добавлены как диагональные по модам: $\gamma_n$ --- декремент амплитуды. При другой нормировке правую часть делят на $N_n=\int_V\rho|\phi_n|^2\dd V$. Если диссипация заметно связывает моды, вместо независимых уравнений получается связанная система.

При гармоническом воздействии $f=\operatorname{Re}[f_\Omega e^{-\ii\Omega t}]$ установившаяся комплексная амплитуда поля равна
\begin{equation}\label{eq:distributedForced}
 u_\Omega(\bm r)=\sum_n
 \frac{\phi_n(\bm r)F_{n,\Omega}}
 {\omega_n^2-\Omega^2-2\ii\gamma_n\Omega}.
\end{equation}
Каждое слагаемое имеет обычный резонансный знаменатель. Вблизи изолированного резонанса может преобладать одна мода, но лишь если её проекция источника не равна нулю. Без потерь точное резонансное возбуждение приводит к росту амплитуды во времени; конечное установившееся значение тогда не существует.

Пространственная структура источника столь же важна, как его частота. Для сосредоточенного воздействия в точке $\bm r_0$ имеем $F_n\propto\phi_n^*(\bm r_0)$: источник в узле не возбуждает соответствующую моду. Распределённый источник, пропорциональный $\rho\phi_m$, возбуждает только моду $m$ в выбранном ортогональном базисе. Это распределённый аналог ортогональности силы нормальному колебанию в~\Qref{1.5.2}.

Для электромагнитных полей резонатора дополнительно соблюдаются уравнения дивергенции: при наличии сторонних зарядов одних поперечных колебательных мод может быть недостаточно, и разложение дополняют потенциальным полем.

\paragraph{Вынужденное напряжение на конце: пример лекции.}
Пусть идеальный источник задаёт $U(0,t)=\mathcal E_0e^{-\ii\Omega t}$, а правый конец линии разомкнут: $U_x(l,t)=0$. Непосредственное решение волнового уравнения даёт
\begin{equation}
 U_\Omega(x)=\mathcal E_0
 \frac{\cos[k(l-x)]}{\cos(kl)},\qquad k=\Omega/v.
\end{equation}
Знаменатель обращается в нуль при $\cos(kl)=0$, то есть на собственных частотах из~\Qref{1.9.1}. В реальной линии потери ограничивают резонансный отклик.

Здесь источник задан \emph{на границе}, а не объёмной правой частью. Поэтому нельзя просто подставить его напряжение вместо $f$ в интеграл проекции. Покажем, как всё же использовать те же собственные функции. Выделим известное граничное значение:
\[
 U(x,t)=\mathcal E(t)+w(x,t).
\]
Тогда $w(0,t)=0$, $w_x(l,t)=0$ --- однородные условия, а уравнение становится
\[
 w_{tt}-v^2w_{xx}=-\ddot{\mathcal E}(t).
\]
Разложив $w=\sum_nq_n(t)\sin(k_nx)$ и используя
$\int_0^l\sin^2(k_nx)\dd x=l/2$,
получим
\begin{equation}
 \ddot q_n+\omega_n^2q_n
 =-\frac{2}{lk_n}\ddot{\mathcal E}(t).
\end{equation}
Для гармонического источника это даёт альтернативное представление того же решения внутри линии:
\begin{equation}
 U_\Omega(x)=\mathcal E_0+
 \sum_{n=1}^{\infty}
 \frac{2\Omega^2\mathcal E_0}
 {lk_n(\omega_n^2-\Omega^2)}\sin(k_nx).
\end{equation}
Выделение граничной части --- самостоятельный вывод из лекционной краевой задачи. Оно показывает смысл модального метода и сохраняет корректное граничное условие. К установившемуся решению отдельно добавляется переходный процесс, определяемый начальными данными.
